\section*{A paid-$U$ stopping theorem with branchwise coordinate reuse}

Let $f:\{-1,1\}^n\to\{-1,1\}$ be Boolean and let $\mathcal T$ be a
finite physical restriction tree.  At an internal node $v$, reached on the
cylinder $E_v$ with probability $p_v$, query a coordinate $i_v$ which is
still unfixed there and write $f_v$ for the restricted Boolean function.
The same physical coordinate may be queried at different incomparable
nodes.  It cannot be queried twice on a single path because it is then
already fixed.

For a Boolean function $g$ and a live coordinate $i$, put
\[
 \delta_iU(g)=U(g)-\frac{U(g_{i=1})+U(g_{i=-1})}{2},\qquad
 q_i(g)=\operatorname{Inf}_i(g),
\]
and let $K_i^{(0)}(g)$ be the empty-upper-set row in the exact identity
\[
 \delta_iU(g)=B_i(g)-\Gamma_i(g)+K_i(g).
\]
Define the non-root residual
\[
 \rho_i(g):=\delta_iU(g)-K_i^{(0)}(g)
 =B_i(g)-\Gamma_i(g)+K_i^{(+)}(g).                 \tag{1}
\]
All terms are evaluated on the genuine Boolean child $g$; no arbitrary
spectral law is introduced.

\paragraph{Exact tree ledgers.}
For the leaf distribution induced by the fair restrictions,
\[
 \sum_{v\in\mathcal T^\circ}p_v\delta_{i_v}U(f_v)
   =U(f)-\sum_{\ell\in\partial\mathcal T}p_\ell U(f_\ell),               \tag{2}
\]
and
\[
 \Pi_I(\mathcal T):=\sum_{v\in\mathcal T^\circ}p_v q_{i_v}(f_v)
   =I(f)-\sum_{\ell\in\partial\mathcal T}p_\ell I(f_\ell).             \tag{3}
\]
The second identity is the real influence decrement; in particular no
queried influence can be charged again at a terminal.

The adaptive Boolean root-row theorem gives, with $c_0=2/\ln2$,
\[
 \sum_{v\in\mathcal T^\circ}p_vK_{i_v}^{(0)}(f_v)
 \le c_0\operatorname{Var}(f).                                         \tag{4}
\]
It already permits reuse of a physical coordinate on incomparable
branches.  Orthogonality of the node wavelets, rather than a bounded
congestion assumption, is what makes the payment single-use.

\paragraph{Theorem (paid boundary and terminal closure).}
Fix absolute $A,\kappa\ge0$ and $0\le b<1$.  Suppose that for every
nonconstant Boolean $f$ there is a finite physical restriction tree such
that
\[
 U(f_\ell)\le\kappa I(f_\ell)\quad\hbox{for every leaf }\ell             \tag{5}
\]
and the one aggregate residual estimate
\[
 \mathcal R(\mathcal T):=
 \sum_{v\in\mathcal T^\circ}p_v\rho_{i_v}(f_v)
 \le A\Pi_I(\mathcal T)+bU(f).                                         \tag{6}
\]
Then, for every Boolean $f$,
\[
 \boxed{U(f)\le
 {\max\{A,\kappa\}I(f)+c_0\operatorname{Var}(f)\over1-b}
 \le {\max\{A,\kappa\}+c_0\over1-b}I(f).}                            \tag{7}
\]
Consequently the reviewed reduction $H\le2I+2U$ gives the uniform FEI
constant
\[
 \boxed{H(f)\le
 \left(2+{2(\max\{A,\kappa\}+2/\ln2)\over1-b}\right)I(f).}            \tag{8}
\]

\paragraph{Proof.}
Add (1) over the internal nodes and use (2), (4), and (6):
\[
 U(f)-\sum_\ell p_\ell U(f_\ell)
 \le c_0\operatorname{Var}(f)+A\Pi_I(\mathcal T)+bU(f).
\]
Now apply (3) and (5):
\[
 (1-b)U(f)
 \le c_0\operatorname{Var}(f)+A I(f)
       +(\kappa-A)\sum_\ell p_\ell I(f_\ell).
\]
Because (3) implies $0\le\sum_\ell p_\ell I(f_\ell)\le I(f)$, the last
two terms are at most $\max\{A,\kappa\}I(f)$.  This proves (7), and (8)
follows from the source reduction.  Constants have $I=U=0$ and are
separate harmless endpoints.

\paragraph{Relation to GU and the exact unpaid line.}
At one node, $G(g)=\sum_i\delta_iU(g)$, while (1) removes only the already
paid root row from the chosen coordinate.  Thus (6) is not a restatement
of GU and does not assume a GU estimate at every intermediate node.  It is
an explicitly sufficient global $U$-drift certificate.  Conversely, a GU
pair $G\le aI+bU$ yields the usual influence-weighted good-coordinate rule
and the source restriction induction $U\le aI/(1-b)$; the present theorem
shows exactly what a stopping proof may replace that local rule by.

The remaining universal obligation is therefore precise: find, for every
Boolean $f$, a finite physical tree satisfying (5)--(6) with one fixed
triple $(A,\kappa,b)$, $b<1$.  Known LTF-terminal modulation trees certify
(5) on a proper family, but do not prove (6) for an arbitrary first
non-flat data collision.  The unresolved summand is the signed residual
$B_i-\Gamma_i+K_i^{(+)}$; no cancellation, terminal term, normalization,
or coordinate-reuse cost has been discarded.
